\chapter*{Anexo: Glosario}
\addcontentsline{toc}{chapter}{Anexo: Glosario}

% TODO: Añadir entradas del glosario con \hypertarget según se vayan
% definiendo términos a lo largo del documento.

\begin{description}

    \item [\hypertarget{bsl}{\textbf{BSL}} (Business Source License)]
    Licencia de software introducida por HashiCorp en 2023 para Terraform. A
    diferencia de las licencias open source, restringe el uso comercial a entidades
    que compitan directamente con el titular de la licencia. Su ambigüedad en la
    definición de «competidor» motivó la creación de OpenTofu como alternativa
    plenamente open source.

    \item [\hypertarget{cicd}{\textbf{CI/CD}} (Continuous Integration / Continuous Delivery)]
    Práctica de desarrollo que automatiza la verificación y entrega del software.
    La integración continua (CI) ejecuta checks de calidad —tests, linting,
    cobertura— en cada cambio integrado al repositorio. La entrega continua (CD)
    automatiza la publicación de artefactos o despliegues cuando el código supera
    esas verificaciones.

    \item [\hypertarget{devcontainer}{\textbf{Devcontainer}}]
    Especificación abierta que define un entorno de desarrollo reproducible mediante
    contenedores. Encapsula las dependencias, herramientas y configuración necesarias
    para trabajar en un proyecto, de forma que cualquier colaborador obtenga un
    entorno idéntico independientemente de su sistema operativo o configuración local.

    \item [\hypertarget{cloud-sprawl}{\textbf{Cloud sprawl}}]
    Proliferación descontrolada de recursos cloud —instancias, bases de datos,
    entornos de desarrollo olvidados— que no están siendo utilizados o monitorizados
    activamente. Genera costes innecesarios y dificulta la gobernanza de la
    infraestructura.

    \item [\hypertarget{finops}{\textbf{FinOps}} (Financial Operations)]
    Disciplina de gestión financiera aplicada al gasto en infraestructura cloud.
    Su objetivo es proporcionar visibilidad sobre los costes, optimizar el uso de
    recursos y alinear el gasto cloud con el valor de negocio generado.

    \item [\hypertarget{iac}{\textbf{IaC}} (Infrastructure as Code)]
    Paradigma de gestión de infraestructura que consiste en definir y provisionar
    recursos mediante ficheros de configuración versionables, en lugar de a través
    de procesos manuales o consolas gráficas.

    \item [\hypertarget{mcp}{\textbf{MCP}} (Model Context Protocol)]
    Protocolo abierto desarrollado por Anthropic que estandariza la forma en que
    los modelos de lenguaje se conectan con herramientas y fuentes de datos externas.
    Permite que un agente de IA invoque funcionalidades de aplicaciones externas
    de forma estructurada y segura.

    \item [\hypertarget{mapper}{\textbf{Mapper}}]
    Componente interno de CloudPorter responsable de resolver las equivalencias entre
    los recursos genéricos definidos en el CloudPorter Manifest y los servicios
    específicos de cada proveedor cloud. Opera de forma determinista: dado un
    recurso y un proveedor, produce siempre la misma traducción.

    \item [\hypertarget{paas}{\textbf{PaaS}} (Platform as a Service)]
    Modelo de servicio cloud en el que el proveedor gestiona la infraestructura
    subyacente y expone al usuario una plataforma para desplegar aplicaciones sin
    necesidad de configurar ni administrar servidores, redes o sistemas operativos.

    \item [\hypertarget{cloudporter-manifest}{\textbf{CloudPorter Manifest}}]
    Fichero en formato YAML que describe la infraestructura de un proyecto de
    forma agnóstica al proveedor cloud. Define recursos de cómputo, red y base de
    datos en términos genéricos que CloudPorter traduce a las plantillas específicas
    de cada plataforma soportada.

    \item [\hypertarget{rag}{\textbf{RAG}} (Retrieval-Augmented Generation)]
    Técnica de inteligencia artificial que combina la generación de texto mediante
    modelos de lenguaje con la recuperación de información de una base de
    conocimiento externa. Permite que el modelo produzca respuestas contextualizadas
    con información específica del dominio, más allá de lo aprendido durante su
    entrenamiento.

    \item [\hypertarget{shared-responsibility}{\textbf{Shared Responsibility Model}}]
    Modelo de seguridad cloud en el que las responsabilidades se dividen entre el
    proveedor y el usuario. El proveedor garantiza la seguridad de la infraestructura
    subyacente; el usuario es responsable de la seguridad de lo que despliega sobre
    ella —configuración de permisos, cifrado de datos, exposición de puertos.

    \item [\hypertarget{skill}{\textbf{Skill}}]
    Comando personalizado de \textbf{Claude Code} que encapsula un flujo de trabajo
    o un conjunto de instrucciones de comportamiento. Las skills se invocan mediante
    una barra oblicua seguida del nombre del comando (\texttt{/nombre}) y permiten
    extender el asistente con comportamientos específicos del proyecto —como el
    estilo de redacción académica o la compilación del documento.

\end{description}

\endinput
